%% CPP 2027 Submission — 0pat Exercise XXXI: Observing the Real and Imaginary Axes through the Euler Circles
\documentclass[sigplan,10pt,anonymous,review]{acmart}
\settopmatter{printfolios=true,printacmref=false}

\AtBeginDocument{%
  \providecommand\BibTeX{{\normalfont B\scshape ib\TeX}}}

\settopmatter{printacmref=false}
\renewcommand\footnotetextcopyrightpermission[1]{}
\pagestyle{plain}

\usepackage{microtype}
\usepackage{booktabs}
\usepackage{amsmath}
% XeTeX (tectonic): acmart loads newtxmath/libertine which already defines \Bbbk;
% reset it so amssymb's definition is accepted.
\let\Bbbk\relax
\usepackage{amssymb}

\begin{document}

\title{Observing the Real and Imaginary Axes through the Euler\\
Circles (Exercise XXXI)}

\author{Anonymous}
\affiliation{%
  \institution{Anonymous}
}
\email{anonymous@example.org}

\begin{abstract}
We observe the real and imaginary axes through the orthogonal Euler
circle frame --- the radius-1 and radius-$i$ circles of the split-complex
numbers --- by solving the intersection equations, not by appealing to
definitions. The observation yields five machine-checked theorems: the
real axis meets the imaginary axis at the single point $(0,0)$
(solving $\langle a, 0\rangle = \langle 0, b\rangle$ gives $a = 0$ and
$b = 0$); the real axis meets the radius-1 circle at $\{\pm 1\}$
(solving $a^2 = 1$); the real axis does not meet the radius-$i$ circle
($a^2 = -1$ has no real solution); the imaginary axis does not meet the
radius-1 circle ($-b^2 = 1$ has no real solution); and the imaginary
axis meets the radius-$i$ circle at $\{\pm i\}$ (solving $-b^2 = -1$).
The observed conclusion: the two axes intersect at $(0,0)$, and
$i = (0,1)$ is the intersection of the imaginary axis with the
radius-$i$ circle --- its unit --- not the intersection of the two axes.
Zero errors, zero \texttt{sorry}.
\end{abstract}

\keywords{Lean, formalization, split-complex, axes observation, intersections, 0pat}

\maketitle

\section{The observation}

The orthogonal Euler circle frame (Exercises XXVII--XXX) consists of the
radius-1 circle $a^2 - b^2 = 1$ and the radius-$i$ circle
$a^2 - b^2 = -1$ in the split-complex plane $a + bj$, $j^2 = +1$. This
exercise observes the real axis ($\mathrm{im} = 0$) and the imaginary
axis ($\mathrm{re} = 0$) through this frame. The method is observational
in the strict sense: every result is obtained by solving the relevant
equation, not by invoking a definition.

\section{The observed table}

\begin{center}
\begin{tabular}{lll}
\toprule
Observation & Result & Equation solved \\
\midrule
real axis $\cap$ imaginary axis & $\{(0,0)\}$ & $\langle a,0\rangle = \langle 0,b\rangle \Rightarrow a=0 \wedge b=0$ \\
real axis $\cap$ radius-1 circle & $\{\pm 1\}$ & $a^2 = 1$ \\
real axis $\cap$ radius-$i$ circle & $\emptyset$ & $a^2 = -1$ (no real solution) \\
imaginary axis $\cap$ radius-1 circle & $\emptyset$ & $-b^2 = 1$ (no real solution) \\
imaginary axis $\cap$ radius-$i$ circle & $\{\pm i\}$ & $-b^2 = -1 \Rightarrow b = \pm 1$ \\
\bottomrule
\end{tabular}
\end{center}

\paragraph{The axes meet at $(0,0)$.} A point on both axes has the form
$\langle a, 0\rangle$ (real axis) and $\langle 0, b\rangle$ (imaginary
axis); equality forces $a = 0$ and $b = 0$. The intersection is the
single point $(0,0)$.

\paragraph{The axes meet the circles in their units.} The real axis
meets the radius-1 circle exactly at its unit $\pm 1$; the imaginary
axis meets the radius-$i$ circle exactly at its unit $\pm i$. In both
cases the axis contributes one equation ($a^2 = 1$, $-b^2 = -1$) and
the circle's unit is the solution.

\paragraph{The cross intersections are empty.} The real axis does not
meet the radius-$i$ circle and the imaginary axis does not meet the
radius-1 circle: both equations ($a^2 = -1$, $-b^2 = 1$) have no real
solution. The units stay on their own axes.

\section{The observed conclusion}

The observation answers the positional question directly: the real and
imaginary axes intersect at $(0,0)$ --- the point obtained by solving
their equations, not by definition. The point $i = (0,1)$ is the
intersection of the imaginary axis with the radius-$i$ circle: it is
the unit of that circle, observed on its own axis. It is not the
intersection of the two axes. The frame's origin $(0,0)$, its units
$\{\pm 1\}$ on the real axis and $\{\pm i\}$ on the imaginary axis, and
the empty cross-intersections form a complete, machine-checked picture
of how the axes sit in the Euler circle frame.

\section{Honest boundary and position}

The content is elementary algebra (linear and quadratic equations over
$\mathbb{R}$); the value is observational. The result concerns the
frame and its axes; no claim is made about the Riemann zeta function.
The frame observes numbers and circles; the analysis of $\zeta$ remains
the open span of the bridge.

\section{Formalization}

The proof is a single Lean file (\texttt{proof.lean}, 190 lines),
importing only \texttt{Mathlib.Data.Real.Basic} and basic tactics. The
split-complex structure (30 lines, as in Exercises XXVII--XXX) is
self-contained. Five theorems, compiled with Lean 4.32.2 / mathlib
v4.32.2, zero errors, zero \texttt{sorry}.

\begin{center}
\begin{tabular}{lll}
\toprule
Theorem & Statement & Proof core \\
\midrule
\texttt{realAxis\_inter\_imaginaryAxis} & axes meet at $\{(0,0)\}$ & ext; simp \\
\texttt{realAxis\_inter\_unitCircle} & real axis meets circle$_1$ at $\{\pm 1\}$ & sq\_eq\_sq; ext \\
\texttt{realAxis\_inter\_imaginaryCircle} & no real-axis point on circle$_i$ & nlinarith \\
\texttt{imaginaryAxis\_inter\_unitCircle} & no imaginary-axis point on circle$_1$ & nlinarith \\
\texttt{imaginaryAxis\_inter\_imaginaryCircle} & imaginary axis meets circle$_i$ at $\{\pm i\}$ & sq\_eq\_sq; ext \\
\bottomrule
\end{tabular}
\end{center}

\section{Conclusion}

The observation is complete: the real and imaginary axes, seen through
the Euler circles, intersect at $(0,0)$; their units $\pm 1$ and $\pm i$
lie on their own axes and their own circles; the cross-intersections are
empty. Every entry of the table is the solution of an equation,
machine-checked. The axes are observed; the analysis of $\zeta$ remains
the open span of the bridge.

\bibliographystyle{ACM-Reference-Format}
\bibliography{refs}

\end{document}
